% ==========================================================================
% GENERATED by scripts/split_paper.py. Do not edit this file.
%
% Source: imported_papers/Llama2_benchmark_sds/llm_at_home.tex
% Part 2: Survey of key technologies
%
% Edits belong upstream, in the published paper. To change what the
% chapter includes, change imports/llmbench.toml and re-run
% the splitter.
% ==========================================================================
\section{Survey of key technologies}
\label{sec:ow}
In this section, we present an overview of the innovations that have enabled model fine-tuning, inference, and performance evaluation of LLMs on a moderate compute budget. We omit to survey how LLMs should be trained from scratch as this undertaking remains extremely costly, being out of scope for most organizations and practical applications. For general design patterns of efficient deep learning, see \cite{menghani2023efficient, tuggener2020design}.

\subsection{Efficient LLM Execution and Adaptation}
\subsubsection{Quantization}
\label{sec:quant}
% TODO Tugg
The biggest hurdle to the wide adoption of LLMs is their VRAM requirements. The parameters and activation states of these models are usually stored in a $32$- or sometimes $16$-bit format. It has been shown that quantization methods that can properly represent the characteristic emergent structures in transformer \cite{vaswani_attention_2017} weights can reduce the parameter and activation formats to an $8$-bit representation with minimal loss of arithmetic accuracy \cite{bnb}. Recently, a $4$-bit NormalFloat data type has been introduced that is theoretically optimal for storing normally distributed weights \cite{qlora}.

When quantizing a fully pre-trained transformer model, the quantization can take the explicit structure of the model weights into account, making quantization even more efficient: GPTQ \cite{gptq} leverages approximate second-order information to generate highly accurate approximations with $3$ to $4$ bits per weight. Activation-aware Weight Quantization (AWQ) \cite{awq} achieves a decrease in quantization error by deriving the importance of each weight through observing their activations before quantization, then protecting the $1\%$ most important weights from being quantized. 

When also training a model (not only running inference), the required VRAM is drastically increased due to the variables required by the optimizer. A straightforward way to address this is to also quantize the optimizer's weights \cite{bnb_opt}.
Quantization is very widely adopted and considered to be the most effective method for reducing VRAM consumption.



% \subsection{Attention Mechanism}

% Attention mechanisms, first introduced for machine translation \cite{bahdanau2014neural}, have recently emerged as pivotal components of neural architectures, and enabled state-of-the-art performance in a remarkably large number of natural language processing tasks (see \cite{attention_review} for a review). Attention mechanisms in language models work by enabling the model to focus on specific parts of the input sequence depending on the need, thereby mimicking human-like selective perception. Effectively, models give varying degrees of importance to different elements in the input sequence dynamically, depending on their relevance to the current step in the sequence generation process. This mechanism comes at a cost of important computational complexity. ``Self Attention,'' which is currently the most commonly used attention mechanism as it is the heart of the transformer architecture \cite{vaswani_attention_2017} (a.k.a. Fast Weight Programmers \cite{DBLP:conf/icml/SchlagIS21}) scales quadratically to sequence length ($O(n^2 d)$), with $n$ sequence length, and $d$ dimension of the embedding). This scaling makes training, and computing inference on large language models computationally very expensive. 

\subsubsection{Low-Rank Approximators}\label{sec:lora}
While quantization methods are often sufficient to run LLM inference on constrained hardware, VRAM constraints prevent fine-tuning LLMs on a single GPU even when using optimizer quantization. Low-rank adaptation of large language models (LoRA) \cite{lora} offers an alternative: LoRA freezes the pre-trained model weights and injects trainable low-rank decomposition matrices into each layer of the Transformer. This way the base network can be run in ``inference mode'' and only the low-rank matrices are trained as a differential update to the base network. AdaLoRa \cite{adalora} develops this concept further and assigns parameter budgets to different weight matrices according to an importance metric. LoKr \cite{LoKr} recognizes limited representation power of above low-rank representations and suggests using the Kronecker product instead.
Low-rank approximators are the key innovation that democratized LLM fine-tuning. In resource-constrained settings they are mandatory.

% \subsection{Other Measurements}


% \subsubsection{Tensor Parallelism and Model Sharding}

% https://huggingface.co/docs/transformers/v4.15.0/parallelism


\subsubsection{Pruning}

A popular strategy to alleviate the high inference cost of deep neural networks is ``model pruning'' \cite{Lecun_Pruning, Hassibi_Pruning, Han_Pruning}: Starting from a pre-trained network, one can iteratively remove redundant parameters and retrain the model while maintaining a test accuracy comparable to the original full-sized network. While Large Language Models are natural candidates for network pruning techniques, they have received little interest from the community compared to the other popular compression techniques such as quantization (see \ref{sec:quant}). A likely explanation for this is that standard pruning methodologies require retraining the network, which is too costly in the case of LLMs. Recent works have however shown promising results with pruning without relying on retraining networks \cite{pruning_llm_1, pruning_llm_2}.

% \subsubsection{Gradient Accumulation}
% Copied from HF: Gradient accumulation is a technique where you can train on bigger batch sizes than your machine would normally be able to fit into memory. This is done by accumulating gradients over several batches, and only stepping the optimizer after a certain number of batches have been performed.

\subsubsection{CPU Parameter Offloading}\label{sec:param_offloading}
CPU offloading in training deep neural networks involves transferring some computations from the GPU to the CPU to alleviate memory constraints and improve training efficiency. Optimal strategies for offloading activations, gradients, and model parameters to the CPU ultimately enable the training of larger models. Microsoft has been at the forefront of this line of research and introduced the open-source library DeepSpeed \cite{deepspeed} as well as ZeRO\cite{zero_offload}. This includes a ``ZeRO-Offload'' feature, which performs optimized CPU parameter offloading, enabling large models with up to $13$ billion parameters to be efficiently trained on a single GPU.
Smart load sharing between CPUs and GPUs is very important in large multi-GPU or even multi-cluster operations. When working with a single GPU it is usually preferable to compute everything on that device to avoid costly data transfers.

\subsubsection{Stateless Optimizers}
% TODO sage
Numerous prevalent optimization algorithms come with the cost of additional parameters. For instance, Adam \cite{kingma_adam_2017} and AdamW \cite{loshchilov_decoupled_2019} involve the computation of both the moving average of first-order moments (mean of gradients) and second-order moments (uncentered variance of gradients) for each parameter. Consequently, they incur the cost of tripling the parameter count. An alternative to quantizing the optimizer weights (see Sec. \ref{sec:quant}) is to use a stateless optimizer like stochastic gradient descent. The drawback of plain SGD is its typically inferior convergence properties. However, this limitation can be mitigated by introducing a learning rate scheduler \cite{loshchilov_sgdr_2017} that systematically reduces the learning rate, ensuring the attainment of a local minimum.
Running a stateless optimizer is an efficient way of reducing VRAM requirements. The incurred loss in convergence speed depends on the problem at hand and it is best to investigate in a preliminary study if it is a worthwhile tradeoff. 

%In this work, we employ AdamW due to its efficiency, notwithstanding the accompanying parameter increase. Our observations indicate that our model can undergo training on a single GPU despite the additional memory overhead incurred by this optimizer. Nonetheless, adopting a stateless optimizer may prove advantageous in alternative scenarios characterized by more constrained memory resources or when training an even larger model.

\subsubsection{Gradient Checkpointing}
% TODO sage
Conventionally, intermediate activations from the forward pass are stored in the VRAM for gradient computation during backpropagation, but this becomes impractical in scenarios with limited memory or very large models. Gradient (or: activation) checkpointing \cite{chen_training_2016, gruslys_memory-efficient_2016} addresses this challenge by selectively caching a subset of crucial intermediate activations for gradient computation, dynamically recomputing the remaining activations during the backward pass. This turns out to make a better trade-off than re-computing all activations: The flexibility of gradient checkpointing in choosing activations for checkpointing allows tailored optimization based on a model's specific memory and computational demands.
Gradient checkpointing can help to further reduce memory consumption when reducing the batch size is no longer possible. It comes, however, at a steep cost in training speed.

%In this work, we do not use gradient checkpointing as our model fits in the memory, and we use a rather small batch size. However, gradient checkpointing would allow the use of larger models and batch sizes at the cost of more inefficient training.

\subsubsection{Gradient Accumulation}
Gradient accumulation \cite{gugger2022accelerate} is a technique that allows training on large ``virtual batches'' that are bigger than the maximum batch size fitting into the available VRAM. This is achieved by accumulating the gradients over multiple batches and only running an optimizer step once enough gradients for the desired ``virtual batch size'' have been collected. 
This is often necessary when working with a limited VRAM budget. Gradient accumulation ensures that the necessary batch size can be reached which reduces the noise in the gradient updates enough for the LLMs to properly converge.

\subsubsection{Software Engineering}
% TODO sage
LLMs typically operate in an autoregressive fashion \cite{vaswani_attention_2017}, predicting the succeeding token based on the prior sequence. Consequently, output generation time increases as the predicted sequence's length grows. In the context of token streaming (as, e.g., in ChatGPT's user interface), the output can be exhibited token by token, as the already produced sequence is directly presented to the user. Furthermore, the autoregressive approach also facilitates continuous batching \cite{yu_orca_2022}. In this method, incoming requests are accumulated in a queue, and after each autoregressive step, they are incorporated into the currently processed batch, provided the memory limit permits. Furthermore, sequences that reach the end-of-sequence token are removed from the batch to free up space, allowing other sequences in the queue to be processed.

These measures are crucial when running LLMs in production; Token streaming should be used when the users are actively waiting for the response, allowing users to promptly read the result, even if it hasn't been fully generated. Continuous batching, on the other hand, is particularly beneficial in scenarios where the timing and quantity of requests are uncertain, for example, in multi-user settings. 

%In this study, we abstain from employing token streaming and continuous batching, instead assessing the model's efficacy with fixed and predetermined batch sizes. Nevertheless, we acknowledge these techniques as crucial when employing LLMs for direct user interactions or multi-user environments.

\subsection{Evaluation}
\label{sec:eval}
% TODO sage
Various metrics such as perplexity \cite{jelinek_perplexitymeasure_1977}, BLEU score \cite{papineni_bleu_2002}, and F1 score have been proposed to gauge the performance of language models on specific tasks, including translation \cite{singh_machine_2017}, and text classification \cite{minaee_deep_2022}.
However, LLM-based assistants have started to exhibit more general capabilities across diverse tasks such as writing, chatting, and coding \cite{anil_palm_2023, touvron_llama_2023, bubeck_sparks_2023, srivastava_beyond_2023}, making evaluating their broad capabilities more challenging.
Existing benchmarks to evaluate LLMs include core-knowledge benchmarks \cite{clark_think_2018, zellers_hellaswag_2019, hendrycks_measuring_2021, sakaguchi_winogrande_2021, chen_evaluating_2021, cobbe_training_2021, zhong_agieval_2023} that typically require LLMs to generate a short, specific answer to benchmark questions that can be automatically validated.
Instruction-following benchmarks \cite{wei_finetuned_2022, mishra_cross-task_2022, wang_super-naturalinstructions_2022, longpre_flan_2023, wang_self-instruct_2023} use more open-ended questions and more diverse tasks to evaluate LLMs after instruction fine-tuning.
Conversational benchmarks \cite{reddy_coqa_2019, feng_mmdialog_2022, kopf_openassistant_2023} focus on dialogue; however, the diversity and complexity of their questions often fall short in challenging the capabilities of the latest models.

Many of these performance metrics primarily center on assessing multi-turn dialogues involving open-ended inquiries, or they fail to pose a sufficient challenge for contemporary cutting-edge models. Although human preferences directly gauge a chatbot's efficacy in multi-turn interactions between humans and AI, this approach lacks scalability. Based on this insight, Zheng et al. \cite{zheng_judging_2023} propose to bridge this gap by using the novel MT-bench benchmark that tests the core capabilities of state-of-the-art models while leveraging an auxiliary LLM model as a judge.
Remarkably, their findings reveal that auxiliary LLM models exhibit a level of agreement comparable to that of human annotators, rendering them suitable for evaluating other LLMs.

In the next section, we apply MT-bench in conjunction with an auxiliary (LLM) to assess the efficacy of a model across diverse tasks. Moreover, this metric proves apt for discerning whether model performance drops during fine-tuning when using, for example, quantization methods, as it not only provides a score per category but also allows pair-wise comparison between models.
